\section{Introduction}
\label{sec:introduction}

Quantum chaos is usually inferred from the spectrum.  After resolving exact
symmetries, nearby levels in a nonintegrable many-body Hamiltonian repel, and
longer-range correlations approach the appropriate random-matrix ensemble.
The same spectral viewpoint underlies much of the connection between quantum
chaos and the eigenstate thermalization hypothesis (ETH)
\cite{srednicki1994,dalessio2016}.  Wave-function chaos can also appear as
universal spectral correlations among reduced-density-matrix eigenvalues,
distinct from the energy-level point process \cite{chenludwig2018}.  The
practical starting point is nevertheless usually a nondegenerate sequence of
energies.

An exactly degenerate manifold removes that starting point.  Let a rank-$D$
subspace have energy $E_0$ and remain separated from its complement by a
positive gap.  Within this subspace every level spacing is zero.  Unfolding is
undefined, the spacing ratio is trivial, and the contribution of the manifold
to the spectral form factor is silent about how its states mix: it depends on
$D$ and $E_0$, not on the orientation of the subspace in the ambient Hilbert
space.  This spectral silence does not classify the Hamiltonian as integrable
or chaotic.  Exact equality of energies can follow from symmetry multiplets, a
common kernel of positive constraints, supercharge cohomology, or dependent
commuting stabilizers.  Each mechanism can coexist with a large and highly
structured complement.

The missing information is carried by the motion of the spectral projector.
For a smooth family $H(\boldsymbol\lambda)$, denote the protected projector by
$P(\boldsymbol\lambda)$ and its complement by $Q=1-P$.  Kato's formulation of
adiabatic transport makes $P$, rather than a basis chosen inside its range,
the natural object \cite{kato1950}.  The Abelian geometric phase
\cite{berry1984}, its non-Abelian extension for degenerate multiplets
\cite{wilczekzee1984}, and the quantum metric
\cite{provostvallee1980} all arise from derivatives of this projector.  The
rectangular operator
\[
 X_a=Q(\partial_aP)P
\]
is the off-fiber projector response: it records how a coupling displacement
rotates protected states into the positive-energy complement.  Products
$X_a^\dagger X_b$ form the non-Abelian quantum geometric tensor, and their
antisymmetric part gives a curvature acting within the protected fiber.
Unlike raw Berry-connection entries, its spectrum and traced contractions are
gauge invariant.

Geometry has already proved sensitive to chaos in nondegenerate many-body
eigenstates \cite{pandey2020}.  More directly relevant here, non-Abelian
Berry curvature was proposed as a chaos probe for exactly degenerate BPS
microstates, where conventional level statistics cannot resolve dynamics
inside the protected sector \cite{chen2026}.  We ask the corresponding
condensed-matter question: what statistical structure can the projector
response retain when degeneracy is enforced algebraically, rather than
obtained by tuning unresolved levels together?  The aim is not to replace
spectral chaos by a single geometric number.  It is to determine which
features of random-matrix behavior survive in the geometry, which remain
sensitive to the parent Hamiltonian, and which exact mechanisms force a
structured negative control.

Two distinctions are essential.  The first separates protection from mixing.
For $h_a=\partial_aH$, the traceless intrafiber perturbation
\[
 S_a=Ph_aP-\frac{\Tr(Ph_aP)}{D}P
\]
tests first-order splitting, whereas $X_a$ tests motion of the subspace.
Along an exactness-preserving transport, $S_a=0$ and the rank and external gap
remain fixed, yet $X_a$ can be nonzero.  Degeneracy and geometric mixing are
therefore independent response channels.  This observation also prevents a
common false inference: a frustration-free or supersymmetric parent need not
have trivial geometry merely because its protected energy is pinned exactly.

The second distinction concerns how a tangent is realized.  An
exactness-preserving transport is a finite family of constraints or generators
for which the common kernel stays exactly degenerate.  In the clustered Hall
parents studied below, differentiating an invertible transport of the constraint
map gives $S_a=0$ identically.  By contrast, the Laughlin fixed-parent
base-point probes apply local potentials at a single exact parent.  They
define the linear response of an isolated rank-$D$ cluster, but a generic
finite displacement splits that cluster.  We keep these two constructions
separate: holonomy and finite-path statements use only the first, while the
second supports only base-point response statistics.

Once the spectrum is silent, the statistical question also divides into
scales.  At the local level we compare curvature eigenvalues with the
finite-rank complex Jacobi process associated with products and compressions
of random subspaces \cite{collins2005}.  Jacobi-like local eigenvalue
repulsion is evidence for correlated mixing within one curvature matrix; it
does not imply that all response entries are independent Gaussian variables.
At the matrix level the two-point covariance and pseudocovariance retain the
tangent labels, complement indices, fiber indices, and parent Hamiltonian.
At fourth order one must subtract all three complex Wick contractions using
that complete covariance.  We refer to the surviving dependence of global
statistics on this tensor structure as global covariance memory.  Local
random-matrix-like repulsion and global covariance memory are compatible, and
their coexistence is one of the central finite-size observations tested here.

Our comparison is organized by the origin of exact degeneracy.  Two-body
Laughlin zero modes are realized in the Kapit--Mueller Chern band and in an
independent continuum lowest-Landau-level construction.  Three-body
Moore--Read zero modes change the clustering constraint and greatly enlarge
the protected rank.  A random cubic supercharge and a local lattice
$\mathcal N=2$ supercharge realize two inequivalent cohomological kernels.
The periodic X-cube stabilizer code supplies an exactly soluble structured
control.  Topological stability under gapped local deformations motivates the
use of projector transport in the Hall setting \cite{hastingswen2005}, but
topology alone does not determine the local curvature statistics.  The models
therefore vary the protection algebra, the locality of the tangent, and the
accessible response channels rather than serving as interchangeable
realizations of one ensemble.

The corrected fourth-order comparison is narrower than the model list. After exact population centering for each registered finite ensemble, channel whitening on a frozen 12-panel training split, and complete three-pairing Wick subtraction on 12 held-out panels, only the registered finite-rank Moore--Read $N=6$ case passes the positive directional gate under a Student-$t_{11}$ reference with a Bonferroni correction over five primary cases. Moore--Read $N=4$ and the lattice-SUSY ranks $D=2,4,8$ do not pass. No corresponding centered full-cumulant statement is assigned to the Laughlin or random-SYK production data. This pattern establishes neither Gaussian closure in the nonpassing cases nor a common ensemble across mechanisms; it fixes the claim boundary of the global comparison.

The resulting claim is deliberately finite.  We find finite-rank,
cross-mechanism evidence that exact spectral degeneracy can coexist with an
informative and, in several stochastic sectors, locally Jacobi-like quantum
geometry.  The response nevertheless remembers its constraint algebra and
operator class through its covariance and branch structure.  The X-cube
control makes the distinction sharp: coefficient deformations have zero
projector curvature, while a chosen unitary transport has deterministic
curvature whose connected variance vanishes.  Thus exact degeneracy alone
does not produce geometric chaos.

These calculations do not constitute an asymptotic Geometric ETH theorem. They neither establish a common scaling limit nor identify a single model/operator class that can be held fixed across all mechanisms. The complete-covariance inference is accepted only for the registered Moore--Read $N=6$ cell whose statistical-unit and familywise gates pass; nonpassing and untested cells are not promoted to conclusions. Section~\ref{sec:protected-response} develops the protected response formalism, Sec.~\ref{sec:statistics} fixes the hierarchy of statistics and inference units, and Sec.~\ref{sec:models} defines the exact mechanisms and tangent classes. The remaining sections present the mechanism-resolved calculations, the topological comparison, and the limits of the geometric interpretation.
